%%============================================================================
%% thesis.tex  --  Adelaide Punt, DPhil thesis (body)
%%
%% Title: Noise and Uncertainty in Molecular Property Prediction:
%%        Robust Models and Adaptive Representations
%%============================================================================

% TODO: REVIEW MODIFICATION (30 Sep) -- tables that ran into the margin now fit: every tabular in a table float is
% wrapped in adjustbox (shrinks only if too wide), the widest headers are split over two lines, and the
% longtable column widths were cut so they fit the 15cm text width.
% TODO: REVIEW MODIFICATION -- every \caption now has a one-line short form in [ ]. Only the short
% form appears in the List of Figures and List of Tables, as in Dablander's thesis.

\chapter{Introduction}\label{chap:intro}
% [notes moved to thesis_notes_moved.md]

% TODO: review, needs manual review
New avenues in the drug discovery pipeline have opened up thanks to advances in large dataset collection, analysis, and computation. Therapeutic targets are identified, potential drug candidates are screened and optimized, and the most promising of these go on to pre-clinical and clinical trials \citep{Lill2007}. There are various methods by which potential drug candidates can be screened. Virtual screening allows for the rapid and large-scale analysis of libraries of compounds, and is a particularly useful approach used to speed up the drug discovery process \citep{Neves2018}. One standard screening method is the prediction of bioactivities and pharmacological properties of such candidates. There is a long history of using quantitative structure--activity relationship (QSAR) models to make these predictions, to act as a filter for compounds with undesirable properties, and thus reduce the time and cost associated with experimental screening \citep{Sabando2021}. QSAR models learn the relationship between a molecule's features, as encoded by a molecular representation, and a property of interest \citep{Sabando2021}. A QSAR model typically consists of a molecular representation paired with a machine learning (ML) model. An ideal molecular representation will allow a model to learn only the relevant information about each molecule with respect to the target property. Representations are either `static' hand-crafted descriptors or circular fingerprints that encode predefined structural and physicochemical features \citep{rogers2010, David2020, sns, avalon}; or learned representations, where the encoding is derived from a predictive model, often taken from the activations of the final layer of a neural network, or a latent space embedding \citep{Sabando2021, kishimoto2023}. 

QSAR models are, of course, sensitive to the quality of the input data on which they are trained. Experimental noise can arise from measurement errors, variability under experimental conditions, varying provenance of the data, and intrinsic biological and chemical variability. Even the very measures of model performance may be flawed if test and validation sets are laden with noise. \citet{Kolmar2021} show that the variance associated with experimental data can contribute more to prediction error than the model’s own error. There is also a limit to how much noise the model can tolerate; \citet{jorner2021} found that once the model error approaches the magnitude of the experimental measurement error on labels, adding more training data yields diminishing returns in predictive performance. In the context of activity prediction, the inherent variability within experimental setups and assays, batch effects, measurement errors, and biological variability introduce additional layers of complexity \citep{Kolmar2021, Fourches2010}. All of this increases the likelihood of encountering noise in the data \citep{Fang2022}. Even the molecular representations in the training data themselves may be incorrect or misleading \citep{Walters2023}, and chemical datasets contain too few repeat measurements of each compound to characterize the distribution the label is drawn from. Thus, models treat a single label as the population mean \citep{Kolmar2021}. Repeat measurements of the same compound-target pair disagree by about $0.68$ log units for pIC$_{50}$, and by about $0.54$ log units for p$K_i$ \citep{Kalliokoski2013, Kramer2012}. Measurements from a single assay lower the mean disagreement from $0.50$ to $0.27$ log units \citep{landrum2024}.
% TODO: the 0.50 baseline sits beside 0.68 (pIC50) and 0.54 (pKi) without saying which comparison it comes from; check landrum2024 and say what 0.50 is measured over.  

% TODO: REVIEW MODIFICATION -- Chapter 1 restructured to match Dablander: prose, then the Outline as Section 1.1.
% What changed, in order:
% - Intro paragraphs 1-2 (QSAR, noise) unchanged.
% - Your paragraph 3 ("This thesis addresses...") and old Section 1.1 ("The questions this thesis asks") are merged into
%   the next paragraph. Your sentence "As we will show, both model and representation choice dominate predictive
%   performance on clean data, however, this relationship inverts..." is cut: it is a Chapter 3 finding, and the
%   Outline states Chapter 3's findings.
% - Your paragraph 4 ("Building on these findings, I began...") is replaced by the two paragraphs from old Section 1.2
%   on the exploratory study and on the distance hypothesis. Your original is kept commented out below.
% - Old Section 1.2's first paragraph (the Chapter 3 variance shares and uncertainty findings) is cut. Kept commented out below.
% - Old Section 1.2's heading is gone. Its notes block moved to thesis_notes_moved.md.
% - The label sec:intro:questions now sits on the questions paragraph; references to it now say Chapter 1.
% - The funding paragraph from "Collaborations and funding" moved to the Acknowledgements in main.tex.

% [TRIM3 2026-09-30] simplified
% CODE INPUT: merged from your paragraph 3 and old Section 1.1, in your wording
% [RESTORE 2026-09-30] Ch1-1 Ch1-2 Ch1-3 Ch1-4
% TODO: review, needs manual review
This thesis addresses the challenge of conducting molecular property prediction when the underlying experimental data contain noisy labels. Most work on noise in QSAR modeling assumes Gaussian label noise and evaluates either representations or model architectures in isolation. \citet{Cortes2015} examined 12 algorithms across multiple noise levels but used fixed representations, and studies that compare representations typically use clean data. No work has systematically compared the relative importance of representation versus architecture choice for noise robustness, asked whether probabilistic methods justify their computational overhead under noise, or tested whether findings generalize across corruption patterns and property types. The first half of this thesis therefore asks four questions.\label{sec:intro:questions} First, how much of the variance in predictive performance do representation and model architecture each explain, and does that division change for the performance retained under increasing label noise? Second, what does the choice of model provide under label noise, and how do deterministic and probabilistic models compare? Third, does the type of label noise matter, or only its quantity? To answer this, we generate and add label noise in several ways, both to clean data and to data with inherent experimental noise. Fourth, what is the impact of artificial label noise on a model's uncertainty estimates? In particular, does the estimated uncertainty rise with increasing label noise, and can it identify which labels are corrupted? Together, these questions aim to identify what, if anything, makes a QSAR model robust to noise. To answer them we introduce NoiseInject, an open-source Python package for benchmarking noise robustness and uncertainty quantification (Section~\ref{sec:noiseinject_pkg}).

% [TRIM3 2026-09-30] simplified
% [TRIM2 2026-09-30] added one sentence on the exploratory LNL study result
% TODO: review, needs manual review
% CODE INPUT: moved from old Section 1.2
Learning with noisy labels (LNL) covers methods for training models when some labels are wrong; noise mitigation, which removes, reweights or corrects labels, is one part of it (Section~\ref{sec:lnl_background}). In an exploratory study carried out earlier in this project, twelve LNL methods were adapted to regression and tested in 221 configurations on the QM9 HOMO--LUMO gap with Gaussian noise, and the 20 configurations that retained the least $R^2$ under noise all used a distance between molecules (Section~\ref{sec:lnl_motivation}).
% - Twelve LNL methods were adapted to regression and tested in 221 configurations on the QM9 HOMO-LUMO gap with Gaussian noise.
% - Several variants judged a label by its neighbors' labels, using a distance between molecules.
% - All 20 configurations with the lowest retention of R^2 under noise used such a distance.
% - This held for each distance tried: Tanimoto on ECFP4, Euclidean on physicochemical descriptors, and Mahalanobis.
% - The methods that held up estimated noise from the model's own loss or uncertainty, not from neighbors.
% - One dataset and one noise type: state it as the motivation, not a general result.
%
% Paragraph 3: why the distance failed, and the step to the hybrid (Chapters~\ref{chap:flaws}, \ref{chap:method} and \ref{chap:similarity})
% - Hypothesis: on a given dataset, most features of a general-purpose representation do not vary or do not relate to the property.
% - Every such feature adds a term to the distance, so irrelevant features dominate it. The growth is linear in the number of features, not exponential.
% - Chapter~\ref{chap:flaws} measures this on base representations and on distances.
% - The proposed fix: select a fixed number of task-relevant features from several representations, and compute distances in that compact space.
% - A first version showed that these selected feature sets also work as ordinary model inputs, with no distance involved.
% - That turned the question into how to choose the features, which is Chapter~\ref{chap:method}.
% - Chapter~\ref{chap:similarity} then returns to distances, noise detection and LNL, including noise mitigation, with the hybrid as the feature space.
% ---- END NOTES ----

% [TRIM3 2026-09-30] simplified
% [TRIM 2026-09-30] shortened
% CODE INPUT: moved from old Section 1.2
% TODO: review, needs manual review
We hypothesized that the fault lay in this distance. On a given dataset, most features of a general-purpose representation either do not vary or do not relate to the property, yet each adds a term to the distance, so together they dominate it (Chapter~\ref{chap:flaws}). We propose instead to compute distances over a hybrid representation: a fixed number of task-relevant features selected from several representations. A first version showed that the selected features also work as ordinary model inputs, so the question became how to choose them (Chapter~\ref{chap:method}).

% [TRIM3 2026-09-30] simplified
% TODO: review, needs manual review
The second half asks whether the representation itself can be adapted to the task:\label{sec:intro:hybrid_questions}
% DECISION FOR ADELAIDE: these five questions do not match the three research questions set out in the Chapter 6 introduction, and questions 2 and 5 have no reported result yet in Chapters 5 or 6 (Chapter 7 answers neither). Align this list with Chapters 5 and 6 once the full runs are in.
\begin{enumerate}
    \item How can the best set of features be selected for a given prediction task?
    \item Do hybrid representations retain more of their predictive performance under label noise than their base representations?
    \item Do the features selected for a hybrid representation help stabilize distance-based noise mitigation?
    \item Does weighting features by their importance improve the construction of hybrid representations, model training and noise mitigation?
    \item When sample removal is suboptimal, can sample weighting or label smoothing handle Gaussian label noise better?
\end{enumerate}

% ---- REMOVED (kept for your review; delete once checked) ----
% This thesis addresses the challenge of conducting molecular property prediction when the underlying experimental data contain noisy labels. The first research question is which combinations of molecular representations and predictive model architectures maintain predictive performance under label corruption, and whether probabilistic methods that estimate uncertainty justify their additional complexity. As we will show, both model and representation choice dominate predictive performance on clean data, however, this relationship inverts under label corruption as predictive architecture becomes the primary factor in determining noise robustness.  Most existing work on noise in QSAR modelings assumes that label noise can be modeled as random error drawn from a Gaussian distribution, and they evaluate either molecular representations or predictive model architectures in isolation. \citet{Cortes2015} examined 12 algorithms across multiple noise levels but used fixed representations. Studies that compare representations are typically done on clean data. No work has systematically compared the relative importance of representation versus architecture choice for noise robustness, assessed whether probabilistic methods justify their computational overhead in noisy settings, or tested whether findings generalize across different corruption patterns and molecular property types.
% 
% Building on these findings, I began exploring techniques involving learning with noisy labels and noise mitigation. Many of these methods incorporate measures of molecular similarity, relying on the assumption that molecules which are similar in structure share similar properties \citep{LopezPerez2024}. However, all of the methods that used simiarity estimates failed, even when I tried varying the molecular representations and vector operations from the standard ECFP4 and Tanimoto coefficient. I hypothesized that to generate accurate similarity estimates, you need a specialized molecular representation who's features relate directly to similarity in that molecular domain. This brings me to my next research question, how do you select the best set of features for a given prediction task? I will outline how I developed an adaptive feature selection methodology that takes in a set of molecular representations, identifies each representation's most important features, and combines a select set of those top features into a new hybrid representation. I will then demonstrate how it improve performance across a range of predictive tasks beyond the small molecule space, and how it can be applied to the context of similarity estimation and downstream noise mitigation methods. 
% 
% The first half of this thesis addresses four questions about the behavior of QSAR models under label noise. First, how does the choice of molecular representation and model architecture explain the variance in predictive performance, and does that division change when you instead consider how much performance a QSAR model is able to retain under increasing label noise? Second, what does the choice of model provide with label noise? Within this question we compare the effects of deterministic and probabilistic modeling. Third, does the type of label noise matter, or just the quantity of it? We explore different techniques of generating and adding label noise to both clean data, as well as data with inherent experimental noise. Fourth, what is the impact of artificial label noise on a model's uncertainty estimation? Does the estimated uncertainty rise with increased label noise and can these uncertainty values point to which labels are corrupted? Together, these questions aim to identify what, if anything, makes a QSAR model robust to noise. Finally, we introduce NoiseInject, an open-source Python package for benchmarking noise robustness and uncertainty quantification in machine learning models.
% 
% % CODE INPUT
% % TODO: Ch3 update -- re-check the Chapter 3 results quoted in this paragraph against the updated chapter.
% The first half of this thesis measures how QSAR models behave as label noise rises (Chapter~\ref{chap:robustness}). We define robustness as the normalized area under the curve of $R^2$ retained as noise increases. Under Gaussian noise on QM9, model architecture explains 49.5\% of the variance in robustness, and representation explains 9.2\%. For $R^2$ at noise level 1.0, where the noise has the same spread as the clean training labels, the two shares are comparable: 29.1\% for architecture and 26.6\% for representation. These shares are measured under noise. \textbf{[RESULT NEEDED: the same variance split on $R^2$ before any noise is added.]} Systematic noise did the most damage: noise shared within a scaffold family, and labels clipped at an assay limit. The shape of the noise distribution barely mattered. Under Gaussian noise on QM9, mean predicted uncertainty rose with noise for every probabilistic model and representation. Per-molecule uncertainty, however, did not reliably pick out which labels were corrupted (Section~\ref{sec:uq}). NGBoost and the Gaussian process fit a separate scale for each molecule. They gave the strongest per-molecule correlation between predicted uncertainty and injected noise. Under censoring, models were most confident on the labels that had been clipped. Uncertainty alone therefore cannot tell us which labels to fix. This led us to methods that judge each label in another way.
% ---- END REMOVED ----


\section{Outline of the Thesis and Contributions}

% ---- NOTES: Outline of the thesis and contributions ----
% Takeaway: one short paragraph per chapter: what it does, what it found (Chapters 3 and 4) or will test (Chapters 5 and 6), and where it is published or planned.
%
% Paragraph 1: Chapter 2 (Background)
% - Introduces, once, what the later chapters rely on: molecular representations, model families, molecular similarity and activity cliffs, noise in molecular data, and the evaluation protocol.
% - Defines the seven noise conditions used in Chapter 3.
% - Ends by introducing NoiseInject, the open-source benchmarking framework released with this work.
%
% Paragraph 2: Chapter 3 (Noise robustness across molecular representations and model architectures)
% - Adapted from the paper "NoiseInject: A Benchmarking Framework to Study Noise Robustness in QSAR Models" by Punt, Hanser, Werner and Morris, in preparation for the Journal of Cheminformatics. Say this in the paragraph.
% - Main task: the HOMO-LUMO gap on QM9. Its labels are computed, so they carry negligible measurement noise.
% - Repeated on three assay datasets with experimental labels: LogD, Caco-2 and hERG Ki.
% - Splits the variance in performance, and in how much performance is retained under noise, between representation and model architecture.
% - Compares seven noise conditions at a matched amount of noise.
% - Compares each model with its probabilistic counterpart, and asks what label noise does to predicted uncertainty.
% - Findings in words only, no numbers:
% - Model architecture matters more than representation for noise robustness. For R^2 under noise, the two matter about equally.
% - The shape of the noise distribution makes little difference. Noise shared within a scaffold family, and labels clipped at an assay limit (censoring), do the most damage.
% - NGBoost and random forests were the most robust models on QM9. Random forests were the most consistent across all four datasets.
% - Predicted uncertainty rises with noise on average, but does not pick out which labels were corrupted.
%
% Paragraph 3: Chapter 4 (The flaws of the similarity assumption)
% - A motivation chapter. It introduces no method, and its measurements are smaller than Chapter 3's.
% - First argument: on a given dataset, standard representations carry many features that do not vary, and some map different molecules to the same vector.
% - Second argument: similarity computed over those representations inherits both problems.
% - Third argument: methods that rely on that similarity fail as a result.
% - These include cleaning labels by comparing each molecule's label with its nearest neighbors' labels, a noise mitigation method tested in the exploratory study of learning with noisy labels. Say plainly that this study comes from the confirmation report.
% - Ends by stating what a representation would need to fix this. That is what Chapter 5 builds.
%
% Paragraph 4: Chapter 5 (Building adaptive hybrid representations)
% - Planned as the first of two method papers. Intended for submission; venue to be confirmed.
% - The method: from several base representations, it selects each one's most important features for the task.
% - It splits a fixed budget of features between the representations and combines the selected features into one representation.
% - Also describes the datasets and the experimental setup.
% - Will test when the hybrid helps, against the best single representation and against an ensemble of one model per representation.
% - Will test why it helps, against PCA, mRMR, recursive feature elimination and random features at the same number of features.
% - Will test whether the selected features can be read against the chemistry, and how stable the selection is across folds.
% - Will test whether search settings found on one dataset carry to another.
% - Will test whether it extends beyond small molecules: antibodies in most depth, and reactions, nucleotide sequences and cancer genomics as proofs of concept.
% - Will test whether it extends to data held at separate sites that cannot be pooled.
% - Result needed: one sentence of findings for Chapter 5.
%
% Paragraph 5: Chapter 6 (Hybrid representations for similarity and label noise)
% - Planned as the second method paper. Intended for submission; venue to be confirmed.
% - Returns to the failures in Chapter 4 and asks whether a hybrid representation repairs them.
% - Similarity: hybrids selected for predicting each molecule's activity from its five nearest neighbors, compared with standard representations.
% - Noise detection and estimation: whether neighbors in a hybrid pick out corrupted labels better than neighbors in a standard representation.
% - Learning with noisy labels, including noise mitigation, one noise type at a time.
% - Result needed: one sentence of findings for Chapter 6.
%
% Paragraph 6: Chapter 7 (Discussion and future directions)
% - Summarizes each chapter.
% - Closing argument: model the noise before choosing a method. No single method fits every dataset, and simple methods are often as good.
% - Then limitations and future directions.
%
% Paragraph 7: Contributions (short: publications and software)
% - Publications: the Chapter 3 paper (in preparation) and the planned papers from Chapters 5 and 6. Point to the Research Papers list in the Summary of Research Contributions rather than repeating it.
% - Software: NoiseInject (open source), and the KIRBy software repository that implements the hybrid search and its experiments.
% - Repositories, licenses and data: point to Additional file~\ref{app:availability}.
% ---- END NOTES ----

% [TRIM3 2026-09-30] simplified
% CODE INPUT
% TODO: review, needs manual review
Chapter~\ref{chap:background} introduces molecular representations, model families, molecular similarity and activity cliffs, and the sources of noise in molecular data. It defines the seven noise conditions used in Chapter~\ref{chap:robustness} and the methods for LNL, and it closes with NoiseInject and the evaluation metrics.

% [TRIM3 2026-09-30] simplified
% CODE INPUT
% [TRIM4 2026-09-30] cut the sentence naming the most robust models; the Chapter 3 Conclusions carry it
% TODO: Ch3 update -- re-check the Chapter 3 results quoted in this paragraph against the updated chapter.
% (30 Sep: findings below rewritten to follow the new Chapter 3 Conclusions.)
% TODO: review, needs manual review
Chapter~\ref{chap:robustness} studies the noise robustness of QSAR models, mainly on the gap between the highest occupied and lowest unoccupied molecular orbitals (HOMO--LUMO gap) in QM9, whose computed labels carry negligible measurement noise, and on three assay datasets: LogD, Caco-2 and hERG $K_i$. It addresses the four questions above, comparing seven noise conditions, all but censoring at a matched amount of noise, and each model with its probabilistic counterpart. Model family explains the largest share of the variance in noise robustness, while representation explains at most 11.2\% of it and a larger share of the variance in predictive performance than of the variance in robustness. The shape of the noise distribution makes little difference, whereas an offset shared within each scaffold family and labels clipped at an assay limit (censoring), the two conditions in which the error is systematic, give the lowest robustness. On the assay datasets, turning a feed-forward network into a Bayesian one raised its robustness at a cost in predictive performance on clean labels. Mean predicted uncertainty rises with noise for almost every probabilistic model on QM9 and for most models on the assay datasets, but per-molecule uncertainty does not identify which labels were corrupted. The exception is censoring: the models are most confident on the clipped labels, so dividing each molecule's error by its uncertainty picks them out better than the error alone.

% [TRIM3 2026-09-30] simplified
% CODE INPUT
% [RESTORE 2026-09-30] Ch1-6
% TODO: review, needs manual review
Chapter~\ref{chap:flaws} motivates the later methods. It introduces no method of its own, and its measurements are on a smaller scale than those of Chapter~\ref{chap:robustness}. On a given dataset, standard representations carry many features that do not vary, and some map different molecules to the same vector. Similarity computed over these representations inherits both problems, and methods that rely on it fail as a result. These include a noise mitigation method, tested in the exploratory study, that cleans each label by comparing it with its nearest neighbors' labels. The chapter ends with what a representation would need to fix these problems, which Chapter~\ref{chap:method} builds.

% [TRIM3 2026-09-30] simplified
% [TRIM4 2026-09-30] planned comparisons cut to one sentence listing them
% CODE INPUT
% DECISION FOR ADELAIDE: Chapters 1 and 6 plan two method papers, but the Research Papers list in the Summary of Research Contributions gives one paper covering Chapters 5 and 6; make them agree.
% TODO: review, needs manual review
Chapter~\ref{chap:method}, planned as the first of two method papers, presents the adaptive hybrid representation. From each of several base representations, the method selects the features most important for the task, divides a fixed budget of features between them, and combines the selected features into one representation. The chapter also describes the datasets and setup. It will compare the hybrid with the best single representation, an ensemble of one model per representation, principal component analysis, minimum-redundancy maximum-relevance selection, recursive feature elimination and random features of the same size. It will also test whether the selected features are interpretable and stable, whether search settings transfer between datasets, and whether the method extends to other domains and to data held at separate sites. Its results so far are preliminary. In local runs on hERG $K_i$ with one default configuration, the hybrid had a lower root mean squared error than the best single representation with a support vector machine and with LightGBM on all three folds, and was level with it for random forest. At the same number of features it was ahead of principal component analysis and random features, but recursive feature elimination was ahead of it.
% TODO: replace the preliminary sentence above with the finding from the full runs for Chapter 5.

% [TRIM3 2026-09-30] simplified
% CODE INPUT
% [RESTORE 2026-09-30] Ch1-7
% TODO: review, needs manual review
Chapter~\ref{chap:similarity}, planned as the second method paper, asks whether a hybrid representation repairs the failures described in Chapter~\ref{chap:flaws}. It compares standard and hybrid representations selected for predicting each molecule's activity from its five nearest neighbors, asks whether neighbors in a hybrid representation identify corrupted labels better than neighbors in a standard one, and tests LNL methods, including noise mitigation, one noise type at a time. Its results are also preliminary. On hERG $K_i$, predicting each molecule's activity from its five nearest neighbors was more accurate in a hybrid than in the best single representation, and, on ECFP4, correcting labels clipped at a limit by one shared amount recovered more than half of the lost performance, whereas deleting the clipped molecules gave a higher error than keeping them. Under Gaussian noise on nanobody thermostability, every one of twelve LNL methods scored higher on the hybrid than any of them did on one-hot, the weakest input.
% TODO: replace the preliminary sentences above with the findings from the full runs for Chapter 6.

% [TRIM3 2026-09-30] simplified
% CODE INPUT
% TODO: review, needs manual review
Chapter~\ref{chap:discussion} summarizes each chapter and argues that the noise in a dataset should be modeled before a method is chosen, since no single method fits every dataset and simple methods are often as good as complex ones. It then gives limitations, future directions and conclusions. % TODO: REVIEW MODIFICATION -- Chapter 7 now ends with Conclusions

% [TRIM3 2026-09-30] simplified
% CODE INPUT
% TODO: review, needs manual review
The software contributions are NoiseInject and the KIRBy repository, which implements the search for hybrid representations and the experiments (Additional file~\ref{app:availability}).


% TODO: REVIEW MODIFICATION -- was "Collaborations and funding"; funding moved to the Acknowledgements.
\section{Collaborations}
% [notes moved to thesis_notes_moved.md]

% CODE INPUT
% TODO: REVIEW MODIFICATION -- funding paragraph moved to the Acknowledgements (main.tex).
% This DPhil was funded by the University of Oxford's UK EPSRC Centre for Doctoral Training in Sustainable Approaches to Biomedical Science: Responsible and Reproducible Research, grant EP/S024093/1. It was also supported by Lhasa Limited, a not-for-profit organization.

% [TRIM3 2026-09-30] simplified
% [TRIM2 2026-09-30] dropped the hERG molecule counts
% CODE INPUT
% TODO: review, needs manual review
This work was supervised by Garrett Morris and co-supervised by Thierry Hanser and St\'ephane Werner at Lhasa Limited, which runs the federated learning platform Effiris \citep{Bassani2023}. Lhasa also supplied the hERG data used in Chapters~\ref{chap:flaws} and~\ref{chap:method}; in Chapter~\ref{chap:method}, two of its sets and the public ChEMBL hERG $K_i$ set are treated as data held at separate sites, as in federated learning \citep{Bassani2023, Hanser2023}.

% [TRIM3 2026-09-30] simplified
% CODE INPUT
% TODO: review, needs manual review
The Oxford Protein Informatics Group (OPIG) pointed us to the antibody and protein datasets used in Chapters~\ref{chap:method} and~\ref{chap:similarity}; all datasets are listed in Additional file~\ref{app:datasets}. The collaboration also produced five posts on OPIG's blog, Blopig, listed under Blog Posts in the Summary of Research Contributions.

% [TRIM2 2026-09-30] "my own work" changed to "the author's own work"
% CODE INPUT
% TODO: review, needs manual review
The adaptive hybrid representation, its software and the experiments in Chapters~\ref{chap:flaws} to~\ref{chap:similarity} are the author's own work.



\pagebreak
